\chapter*{Tóm tắt}
\addcontentsline{toc}{chapter}{Tóm tắt}
\changefontsizes[16pt]{13pt}

Hệ thống gợi ý hội thoại (Conversational Recommender System - CRS) đóng vai trò trung tâm trong việc giải quyết bài toán thông tin cá nhân hóa thông qua tương tác ngôn ngữ tự nhiên. Tuy nhiên, sự kết hợp giữa các Mô hình ngôn ngữ lớn (LLM) và các kho tri thức cấu trúc vẫn tồn tại những rào cản về tính xác thực của dữ liệu và năng lực suy luận linh hoạt trong các kịch bản thực tế.

Khóa luận này đề xuất hệ thống \textbf{ARGOS} (Adaptive Reasoning and Graph Orchestration System), một khung làm việc đa tác tử được kế thừa và phát triển từ mô hình tiền đề \textbf{GRACE} (công trình nghiên cứu của chính tác giả đã được công bố tại hội nghị quốc tế SOICT 2025). Khác với các kiến trúc dòng chảy tĩnh vốn dễ gặp bế tắc khi đối mặt với các yêu cầu phức tạp, ARGOS vận hành dựa trên cơ chế cộng tác giữa các tác tử chuyên biệt: Orchestrator (Điều phối), Graph Agent (Suy luận đồ thị), Critic Agent (Kiểm định) và Relaxation Agent (Thỏa hiệp). Kiến trúc này cho phép hệ thống thực thi các vòng lặp phản hồi để tự phát hiện ảo giác, đồng thời tự động điều chỉnh chiến lược truy xuất khi gặp phải các ràng buộc mâu thuẫn từ phía người dùng.

Kết quả đánh giá thực nghiệm trên các bộ dữ liệu tiêu chuẩn ReDial và INSPIRED chứng minh rằng ARGOS không chỉ bảo toàn tính chính xác tuyệt đối nhờ sự ràng buộc vào Đồ thị tri thức, mà còn cải thiện đáng kể khả năng hiểu ý định và tính linh hoạt trong hội thoại so với kiến trúc tiền đề. Nghiên cứu này khẳng định tính hiệu quả của phương pháp tiếp cận đa tác tử trong việc xây dựng các hệ thống gợi ý hội thoại thông minh, tin cậy và có khả năng mở rộng.

\vspace{-0.5cm}
\begin{flushleft}
  \textit{\textbf{Từ khóa:} Hệ thống gợi ý hội thoại, Đồ thị tri thức, Mô hình ngôn ngữ lớn, Hệ thống đa tác tử (Multi-Agent System), GRACE, ARGOS.}

\end{flushleft}
